\documentclass[12pt, a4paper]{letter}
\usepackage[utf8]{inputenc}
\usepackage[T1]{fontenc}
\usepackage{geometry}
\usepackage{amssymb}
\usepackage{hyperref}

\geometry{margin=1in}

\hypersetup{colorlinks=true, linkcolor=blue, urlcolor=blue}

\signature{Rafael D. De Paz \\ Independent Researcher \\ \texttt{hi@logos.pub}}
\address{Rafael D. De Paz \\ Independent Researcher \\ \url{https://logos.pub}}

\begin{document}
\begin{letter}{Editorial Board \\ Foundations of Physics \\ Springer}

\opening{Dear Editors,}

I am writing to submit the enclosed manuscript, entitled

\begin{center}
\textbf{``Information-Theoretic Gravity: \\
Proton Mass from Vacuum Geometry and the \\
Emergence of the Schwarzschild Metric,''}
\end{center}

\noindent for consideration for publication in \textit{Foundations of Physics}.

\medskip

\textbf{Summary.} The paper presents a derivation of gravity as an emergent
informational phenomenon. Using the holographic mass framework of Haramein,
Guermonprez, and Alirol (2023), we derive the proton rest mass from the
quantum vacuum energy density through a two-surface screening mechanism
containing \emph{zero free parameters}. The predicted proton charge radius
($r_p = 0.84132$ fm) agrees with CODATA 2018 to better than $0.01\%$.
We prove that the Schwarzschild metric emerges naturally from the first
screening surface and provide an independent topological derivation via
Gauss-Codazzi embedding in $S^3$.

\medskip

\textbf{Suggested Reviewers.}
\begin{enumerate}
  \item \textbf{Prof. Erik Verlinde} (University of Amsterdam) ---
    Author of the entropic gravity framework.
  \item \textbf{Prof. Gerard 't Hooft} (Utrecht University) ---
    Pioneer of the holographic principle.
  \item \textbf{Prof. Leonard Susskind} (Stanford University) ---
    Co-developer of the holographic principle and string landscape.
\end{enumerate}

\medskip

This manuscript has not been previously published and is not under consideration
at any other journal. I confirm no conflicts of interest.

\closing{Respectfully submitted,}
\end{letter}
\end{document}
